\section*{Overview}

Mo's algorithm is an offline range-query technique for static arrays. It does not invent a new query formula. It
reorders the queries so one maintained interval can be reused efficiently from one query to the next.

\section*{Problem-Driven Motivation}

This technique appears when a problem has:

\begin{itemize}
  \item a static array,
  \item many subarray queries,
  \item an answer that is awkward for Fenwick trees or segment trees,
  \item but a statistic that is easy to update when one element enters or leaves the current interval.
\end{itemize}

The classic example is distinct values in subarrays. A segment tree does not naturally support ``how many values have
positive frequency'' with easy merges. But if I already know the current interval, adding or removing one endpoint is
trivial.

\section*{Recognition Pattern}

Mo's algorithm is usually worth considering when:

\begin{itemize}
  \item the array is static or nearly static,
  \item offline processing is allowed,
  \item the answer can be maintained by symmetric \texttt{add(index)} and \texttt{remove(index)} operations,
  \item \(O((n+q)\sqrt{n})\) is acceptable and a cleaner online \(\log n\) structure is not obvious.
\end{itemize}

Typical signals are:

\begin{itemize}
  \item distinct values,
  \item frequency-score queries,
  \item range statistics with local add/remove logic,
  \item tree versions after Euler tour reduction.
\end{itemize}

\section*{Derivation}

If I answer queries in arbitrary order, the maintained interval may jump wildly, so the total number of endpoint moves is
too large. The only real freedom in an offline problem is the order of the queries, so Mo's algorithm spends that freedom
to reduce pointer movement.

Split indices into blocks of size about \(\sqrt{n}\). Sort queries by:

\begin{itemize}
  \item block of the left endpoint,
  \item then by right endpoint, often alternating direction per block to reduce movement in practice.
\end{itemize}

After sorting, maintain one active interval \([curL, curR]\). To reach the next query \([l,r]\), move the endpoints one
step at a time and call only:

\begin{itemize}
  \item \texttt{add(index)},
  \item \texttt{remove(index)}.
\end{itemize}

The entire method depends on one invariant:

\begin{quote}
After every single pointer move, the maintained state matches the current active interval exactly.
\end{quote}

\section*{Worked Problem}

\paragraph{Problem.}
For each query \([l,r]\), count how many distinct values appear in \(a[l \dots r]\).

\paragraph{Why naive fails.}
Scanning each interval independently is \(O(nq)\), which is too slow at \(n,q \le 2 \cdot 10^5\).

\paragraph{Why prefix methods are awkward.}
Distinct-count is not an additive statistic. There is no simple formula that subtracts two prefix answers.

\paragraph{Mo invariant.}
Maintain:

\begin{itemize}
  \item a frequency array \(\texttt{freq[value]}\),
  \item one integer \(\texttt{distinct}\).
\end{itemize}

Then:

\begin{itemize}
  \item on add of value \(x\), if \(\texttt{freq[x]}\) becomes \(1\), increment \(\texttt{distinct}\),
  \item on remove of value \(x\), if \(\texttt{freq[x]}\) becomes \(0\), decrement \(\texttt{distinct}\).
\end{itemize}

\paragraph{Final algorithm.}

\begin{itemize}
  \item coordinate-compress values if needed,
  \item sort queries in Mo order,
  \item move the current interval toward each query,
  \item record \(\texttt{distinct}\) when the interval matches.
\end{itemize}

\section*{Implementation Reasoning}

The logic of \texttt{add} and \texttt{remove} matters more than the sorting comparator. If those two functions are not
perfect inverses of each other, Mo's algorithm quietly returns nonsense.

Practical implementation choices:

\begin{itemize}
  \item compress large values first, or the frequency array becomes impossible to allocate,
  \item use alternating right-endpoint order inside each block for better constants,
  \item keep intervals consistently inclusive or half-open; mixing the two is a common bug source.
\end{itemize}

The code below is intentionally the classic distinct-values version because it demonstrates the invariant cleanly.

\section*{Correctness Intuition}

Reordering queries does not change their meaning. It only changes how much work is reused between them. As long as each
endpoint move updates the maintained state correctly, the final state after all moves equals the answer for the new
interval.

\section*{Complexity Analysis}

With block size about \(\sqrt{n}\) and \(O(1)\) add/remove:

\begin{itemize}
  \item sorting: \(O(q \log q)\),
  \item total interval movement: about \(O((n+q)\sqrt{n})\),
  \item memory: usually \(O(n)\) or \(O(\text{value range after compression})\).
\end{itemize}

\section*{Common Pitfalls}

\begin{itemize}
  \item Forgetting that the whole method is offline.
  \item Writing \texttt{add} and \texttt{remove} that are not exact inverses.
  \item Using raw values without compression.
  \item Choosing Mo's algorithm when a simpler online structure already solves the problem more directly.
\end{itemize}

\section*{Variants and Failure Modes}

\begin{itemize}
  \item Mo with updates adds a time dimension and is significantly more fragile.
  \item Tree Mo first reduces the tree to an Euler-tour order.
  \item Hilbert-order sorting often improves constants.
  \item If add/remove are not cheap, Mo's ordering alone does not save the solution.
\end{itemize}

\section*{Practice Problems}

\begin{itemize}
  \item Distinct values on subarrays.
  \item Frequency-score range queries.
  \item Offline array or tree queries with local add/remove maintenance.
\end{itemize}

\section*{References}

\begin{itemize}
  \item \href{https://cp-algorithms.com/data_structures/sqrt_decomposition.html}{cp-algorithms: Sqrt Decomposition and Mo's Algorithm}.
  \item \href{https://usaco.guide/plat/sqrt}{USACO Guide: Square Root Decomposition and Mo's Algorithm}.
  \item \href{https://codeforces.com/catalog}{Codeforces Catalog}.
\end{itemize}
